% =====================================================================
%  GMN3 §8  模空间与全纯函数的例子
%  （Some examples of the moduli spaces and holomorphic functions）
%  原文：Lamination-Paper.tex 4686--5050 行；PDF p.73--78（印刷 p.72--77）
%  编号公式 (8.1)--(8.19) 共 19 条，全部在本块内；
%  跨块引用（节号、式号）暂写原文编号并加 % TODO，合稿时统一改 \ref/\eqref
%  label 名照抄原文：eq:N3-manifold, eq:A1-covering-maps, eq:nf0-matcond,
%  eq:monod-equa-nstar, eq:moduli-equation, eq:CH-thm, eq:A1-covering-maps-nstar
% =====================================================================

\section{模空间与全纯函数的例子}
\label{sec:Examples-of-M}

本节在本文所用的那些具体理论例子中，描述模空间 $\CM$，以及它们与第 6.3、6.4、7.3.1、7.3.2 节% TODO \ref{sec:moreonmoduli}、\ref{sec:twisted-vev}、\ref{subsec:Hitchin-Isogenous}、\ref{sec:twisted-loc}
所引入的空间 $\CM_{\CL_i}$ 与 $\widetilde{\CM}$ 的关系。我们还将描述这些空间上的函数代数，如何与第 4 节% TODO \ref{sec:Formal-line-operators}
所述线算子的非交换代数、以及第 10 节% TODO \ref{sec:Examples}
所述的叶层相关联。

\subsection{$U(1)$ 规范理论与周期 Taub-NUT}

我们从第 4.2 节% TODO \ref{subsec:Formal-line-ops-U1}
的 $U(1)$ 理论开始。$\CM$ 的环面纤维化结构可描述如下。一个耦合了电荷为 $1$ 的超多重态的 $U(1)$ 规范理论，其库仑分支 $\CB$ 在原点 $Z_{0,1}=a=0$ 处有一个奇点，在那里该超多重态是无质量的。它在离原点很远处还有一种非物理的行为，原因在于这个理论不是渐近自由的。在 $\CB$ 中离原点足够近处，理论是有意义的。作维数约化后，四维规范场给出两个圆周值标量：电 Wilson 线 $\tilde\theta_e$ 与磁 Wilson 线 $\tilde\theta_m$。磁圆周在库仑分支的原点处收缩成一点。

$\CM$ 上的度量被「周期 Taub-NUT」空间（又称「Ooguri–Vafa」空间）\cite{Ooguri:1996me,Seiberg:1996ns} 的三维单圈修正抹平。作了这些修正之后可以看出，磁圆周只在余维数为 3 的地方收缩——即在四维库仑分支的原点（$a=0$）处，以及在电圆周上的某个特定点（$\tilde\theta_e=0$）处。

我们提到过的关于不同模空间 $\CM_\CL$ 的那些微妙之处，在这个例子里并不存在：本例中没有值得考虑的对 $\CL$ 的选取。

现在来看 $\CM$ 上的函数代数。在一般的复结构下，该流形可以等同于复流形 $V_+V_-=1-U$，其中 $U=\CY_{0,1}=\exp\left(\frac{\pi R}{\zeta}a+i\tilde\theta_e+\pi R\zeta\bar a\right)$ 取值于 $\IC^*$，而 $V_\pm$ 是更复杂一些的全纯函数，取值于 $\IC$。它们在 $\zeta$ 平面上的渐近性质是：在 $\zeta$ 平面中满足 $\Im(Z_{0,1}/\zeta)>0$ 的扇区里 $V_+=\CY_{1,0}$，而在相对的扇区里 $V_-=\CY_{-1,0}$。我们立刻认出：$V_+^pU^q$ 与 $V_-^pU^q$ 按 $\CY_\gamma$ 的展开、以及它们的环关系，与第 4.2 节% TODO \ref{subsec:Formal-line-ops-U1}
讨论过的形式线算子 $F_{p,q}$ 与 $F_{-p,q}$ 所算出的结果相符——只要取 $y=-1$。

考察一个带电荷为 $Q$ 的超多重态的 $U(1)$ 规范理论所对应的几何，也很有启发性。此时流形的方程取 $V_+V_-=(1+(-U)^Q)^Q$ 的形式，并有一组 $Q$ 个 $A_{Q-1}$ 奇点，位于 $U^Q=(-1)^{Q-1}$：磁圆周在电圆周上的 $Q$ 个位置处收缩，并且它比以前小了 $Q$ 倍，于是在它收缩的地方形成 $A_{Q-1}$ 奇点。注意，这与环关系 (4.20)% TODO \eqref{eq:Q-U1-algebra}
的 $y=-1$ 特殊化漂亮地吻合。（为看出这一点，注意到在 $y=-1$ 时有 $[n]=(-1)^{n-1}n$、$[n]! =(-1)^{\half n(n-1)}n!$，以及 $\genfrac{[}{]}{0pt}{}{n}{j}_{y=-1}=(-1)^{j(n+1)}\binom{n}{j}$。）

\subsection{$N=3$ Argyres--Douglas 理论}\label{subsec:N3-AD-Manfld}

现在转向第 4.3 节% TODO \ref{sec:formaln3}
讨论过的 $N=3$ 的 AD 理论。我们先描述其纤维化结构。

当紧化圆周的半径 $R$ 很大时，这个理论的三维库仑分支有一个直观的结构：库仑分支中两个奇点各自的邻域，都像是 $U(1)$ 理论耦合单个电荷为 $1$ 的粒子时的库仑分支。两个奇点以完全相同的方式被抹平，当然在每个奇点处收缩的是不同的圆周。粗略地说，在一个奇点上方的电圆周的某一点处，磁圆周收缩；而在另一个奇点上方的磁圆周的某一点处，电圆周收缩。Wilson 线的平移对称被完全破坏。与 $U(1)$ 的例子一样，本例中没有值得考虑的对 $\CM_{\CL_i}$ 的选取。

由 \cite{Gaiotto:2009hg} 第 9 节所述的与 Hitchin 系统的对应可知，我们可以把 $\CM$ 描述为球面上一个 $SU(2)$ Hitchin 系统的模空间，该系统只有一个不规则奇点，其阶数高到足以在该奇点周围给出五个斯托克斯扇区。作为复流形，$\CM$ 由方程 $M=1$ 给出，其中单值矩阵 $M$ 分解为
\begin{equation}
M  = \begin{pmatrix} 1 & U_1 \\  0  & 1 \\ \end{pmatrix}\begin{pmatrix} 1 & 0 \\  - U_4  & 1 \\ \end{pmatrix}\begin{pmatrix} 1 & U_2 \\  0  & 1 \\ \end{pmatrix}\begin{pmatrix} 1 & 0 \\  - U_5  & 1 \\ \end{pmatrix}\begin{pmatrix} 1 & U_3 \\  0  & 1 \\ \end{pmatrix}\begin{pmatrix} 0 & 1 \\  - 1  & 0 \\ \end{pmatrix}.
\end{equation}
这里面的三角矩阵是斯托克斯矩阵，而最后一个矩阵是形式单值。$U_i$ 是模空间上的全纯函数。它们的编号是为日后方便而取的。我们还将把编号延拓，定义 $U_{i+5}=U_i$。方程 $M=1$ 稍作整理后\footnote{证明这一点的一种办法，是把三个因子移到方程 $M=1$ 的一侧。这样就得到关于矩阵元的四个方程，容易看出它们正是 $i=2,3,4,5$ 时的方程 \eqref{eq:N3-manifold}。要得到最后一个方程，把 $i=4$ 的方程乘以 $U_2$，并注意只要 $U_3$ 非零，就得到 $i=1$ 的方程。类似地，若把 $i=3$ 的方程乘以 $U_5$ 再除以 $U_4$，就得到 $i=1$ 的方程。但 $U_3$ 与 $U_4$ 不可能同时为零。}，取一个漂亮的形式：
\begin{equation}\label{eq:N3-manifold}
U_{i-1} U_{i+1} = 1+U_i.
\end{equation}

可以证明，这五个方程定义了 $\IC^5$ 中一个光滑的二维复子流形。如果把五个坐标与五边形的五个顶点相对应，那么相邻顶点上所附的坐标绝不会同时为零。有五个特殊的除子 $U_i=0$，它们两两相交，呈五边形构型，交于五个特殊点，即 $U_1=U_3=0$、$U_2=U_4=U_5=-1$，以及它们在指标循环置换下的像。

把方程 (4.29)% TODO \eqref{eq:Deformd-N3-rels}
与 \eqref{eq:N3-manifold} 相比较，可以看出：取 $y\to 1$ 时的极限 $F_i\to U_i$，或取 $y\to -1$ 时的极限 $F_i\to -U_i$，就能还原出交换函数代数。这里 $\CM$ 与 $\tCM$ 似乎是相同的。我们将在第 10.1 节% TODO \ref{subsec:ADN3-Laminations}
中，把 $-U_i$ 按 $\CY_\gamma$ 的展开与 $U_i$ 按 $\tCY_\gamma$ 的展开，同形式线算子逐项对应起来。

\subsection{$SU(2)$ 与 $SO(3)$ 的 $N_f=0$ 规范理论}\label{subsec:SU2-Functions}

现在转向第 4.4 节% TODO \ref{subsec:formal-SU2}
讨论过的 $SU(2)$ 理论。我们先描述 $u$ 平面上的纤维化结构。

平直空间中的 $SU(2)$ 与 $SO(3)$ 的 $N_f=0$ 规范理论，只差在允许的局部线算子集合的选取上。我们已经看到，实际上有三种可能，而第三种可能在把紫外 $\theta$ 角平移 $\pi$ 之后，可以等同于 $SO(3)$ 的 $N_f=0$ 理论。三个理论的三维库仑分支 $\CM_1$、$\CM_2$、$\CM_3$，只差在 Wilson 线参数的精确周期上。我们先描述 $SU(2)$ 情形的模空间 $\CM_1$。这个理论的四维库仑分支 $\CB$ 与 $N=3$ AD 理论的情形惊人地相似，只是在奇点处变为无质量的粒子，其电荷满足 $\langle\gamma_1,\gamma_2\rangle=2$。按惯例，可以取 $\gamma_1=(1,0)$ 与 $\gamma_2=(-1,2)$。所有 BPS 粒子的电电荷都是偶数，因而瞬子修正、从而 $\CM_1$，在 $\IZ_2$ 平移对称 $\tilde\theta_e\to\tilde\theta_e+\pi$ 下不变。由于 $\gamma_i$ 经 $SL(2,\IZ)$ 变换与一个纯电电荷相联系，大半径处的瞬子修正会在每个奇点周围给出标准的周期 Taub-NUT 几何。

对 $SO(3)$ 规范理论，我们则会说 $\gamma_1=(2,0)$ 与 $\gamma_2=(-2,1)$：相对于 $SU(2)$，磁电荷加倍而电电荷减半。于是 $\CM_2$ 有一个不同的 $\tilde\IZ_2$ 平移等距，作用在 $\tilde\theta_m$ 上。由于 $\gamma_1$ 可被 $2$ 整除，大半径处 $\CB$ 的第一个奇点附近的局部几何，预期会有两个 $A_1$ 奇点，它们被 $\tilde\IZ_2$ 置换。对 $\CM_3$，两个奇点的角色互换。

为了把各个 $\CM_i$ 联系起来，我们可以做两件事。一种可能是用它们的 $\IZ_2$ 平移对称作商，把它们的环面纤维归约到共同的形式——由格 $\Gamma$（而非 $\Gamma_{\CL_i}$）构造。所得的流形 $\CM$ 有两个 $A_1$ 奇点，对应于四维库仑分支中的两个奇点。第二种可能是把它们的环面纤维「加倍」，得到更大的环面，以 $\Gamma\otimes\IR$ 中 $\Gamma_{\CL_i}$ 的并为模型。所得的流形 $\CM_{mx}$ 是光滑的，并具有 $\IZ_2\times\tilde\IZ_2$ 平移等距。此时图 (6.12) 变成% TODO \eqref{eq:covering-maps}
\begin{equation}\label{eq:A1-covering-maps}
\begin{matrix}
    &  &  \CM_{mx} &  &   \\
    & \swarrow & \downarrow  &  \searrow & \\
    \CM_{1} &  & \CM_{2} &  & \CM_{3} \\
    & \searrow & \downarrow  &  \swarrow & \\
    &  &  \CM &  &   \\
\end{matrix}.
\end{equation}
于是 $\CM_{mx}$ 对 $\IZ_2$ 作商得到 $\CM_1$，对 $\tilde\IZ_2$ 作商得到 $\CM_2$，对对角 $\IZ_2$ 作商得到 $\CM_3$。

这些流形与预期的 Hitchin 系统解的模空间相比如何？我们马上就会看到，$\CM_1$ 是 $SU(2)$ Hitchin 系统的标准模空间，而 $\CM_{mx}$ 是它的一个自然扩张。这个例子提示：不规则奇点的情形也许比正则奇点的情形要稍微复杂一些——即使 $C$ 的亏格为零，也可能存在规范群与 $\CM_\CL$ 的选取。

该 Hitchin 系统在球面上有两个不规则奇点，每个只有一个斯托克斯扇区。条件 $M_1M_2=1$ 变成
\begin{equation} \label{eq:nf0-matcond}
W \begin{pmatrix} 1 & U_1 \\  0  & 1 \\ \end{pmatrix}\begin{pmatrix} 0 & 1 \\  - 1  & 0 \\ \end{pmatrix} W^{-1} \begin{pmatrix} 1 & U_2 \\  0  & 1 \\ \end{pmatrix}\begin{pmatrix} 0 & 1 \\  - 1  & 0 \\ \end{pmatrix}=1.
\end{equation}
这里 $W$ 是行列式为 $1$ 的矩阵。还剩下一个小小的「规范自由度」$W\to -W$。这个 Hitchin 系统模空间上真正的全纯函数必须在这 $\IZ_2$ 变换下不变，但眼下我们先忽略这一点。$W$ 与 $U_{1,2}$ 的矩阵元可以看成某个略微扩大了的模空间上的全纯函数，这个模空间将与 $\CM_{mx}$ 重合。

把 \eqref{eq:nf0-matcond} 稍作整理，可以看出它蕴涵 $U_1=U_2=:U$ 以及 $\Tr W=0$，即可以写成
\begin{equation}
W = \begin{pmatrix} V_1 & V_0 \\  -V_2  & -V_1 \\ \end{pmatrix}.
\end{equation}
此外 \eqref{eq:nf0-matcond} 还蕴涵 $V_0+UV_1+V_2=0$；又因 $\det W=1$，有 $V_0V_2=1+V_1^2$。二维空间 $\CM_{mx}$ 以 $U,V_0,V_1,V_2$ 为坐标，并受这两个方程约束。

我们可以把 $V_0,V_1,V_2$ 延拓成一族无穷多个有用的函数 $V_n$：
\begin{equation}
W \left(\begin{pmatrix} 1 & U \\  0  & 1 \\ \end{pmatrix}\begin{pmatrix} 0 & 1 \\  - 1  & 0 \\ \end{pmatrix} \right)^n=\begin{pmatrix} V_{n+1} & V_n \\  -V_{n+2}  & -V_{n+1} \\ \end{pmatrix}.
\end{equation}
这些定义的自洽性要求对所有的 $n$ 都有 $V_{n-1}+UV_n+V_{n+1}=0$。此外，由于 $\det W=1$，对所有的 $n$ 都有关系式 $V_{n-1}V_{n+1}=1+V_n^2$。

单值的迹有简单的形式 $\Tr M_1=-U$，它应当对应于基础表示 $SU(2)$ Wilson 圈的期望值。我们有两个自然的对称：一个 $\IZ_2$ 作用为 $V_n\to -V_n$，一个 $\tilde\IZ_2$ 作用为 $U\to -U$、$V_n\to(-1)^nV_n$。前者正是我们为得到 $SU(2)$ Hitchin 系统解的模空间所必须作的商；它作用自由，所以模空间是光滑的。另一方面，$\tilde\IZ_2$ 或 $\IZ_2\times\tilde\IZ_2$ 中的对角子群确实有不动点：或者 $U=0$ 且 $V_{2n}=0$，这意味着 $V_{2n+1}=\pm(-1)^ni$；或者 $U=0$ 且 $V_{2n+1}=0$，这意味着 $V_{2n}=\pm(-1)^ni$。因此，对 $\tilde\IZ_2$ 作商、或对 $\IZ_2\times\tilde\IZ_2$ 中的对角子群作商，得到的空间有两个 $A_1$ 奇点，它们被剩余的 $\IZ_2$ 对称互换。对两个对称都作商，得到的空间有两个互不相关的 $A_1$ 奇点。这与规范理论分析所预期的完全一致：不作任何商的模空间似乎与 $\CM_{mx}$ 重合，$\IZ_2$ 商与 $\CM_1$ 重合，而 $\tilde\IZ_2$ 商与对角商则似乎分别与 $\CM_2$、$\CM_3$ 重合。

现在，把 $\CM_{mx}$ 与 $\CM_1$ 上的函数环同第 4.4 节% TODO \ref{subsec:formal-SU2}
关于形式线算子的结果作比较，是颇有意思的。为得到交换环，可以特殊化 $y=\pm1$。若取 $y=1$，就可以与 $\CM$ 上的函数 $V_n$ 比较。此时当 $y\to 1$ 有 $\hat V_n\to V_n$ 以及 $W\to -U$。把上面的关系式与方程 (4.50)% TODO \eqref{eq:SU2-formring-2}
和 (4.51)% TODO \eqref{eq:SU2-formring-3}
比较。当我们作商走到 $\CM_1$ 时，必须考虑由 $V_n$ 生成的偶函数。此时有 $W\to -U$、$G_{2n+1}\to V_n^2$ 以及 $G_{2n}\to V_{n-1}V_n$。在讨论 $\CM_1$ 上的函数时，也可以在由 $G$ 生成的代数中取 $y=-1$。此时 $G_{2n+1}\to -V_n^2$、$G_{2n+2}\to -V_nV_{n+1}$，而 $W\to U$。确实，有了这一对应，我们可以验证：方程 (4.43)、(4.44)、(4.52)、(4.53)% TODO \eqref{eq:rec-rels-1}、\eqref{eq:rec-rels-2}、\eqref{eq:SU2-formring-4}、\eqref{eq:SU2-formring-5}
在特殊化到 $y=-1$ 后，被相应的经典函数所满足。

\subsection{$SU(2)$ 与 $SO(3)$ 的 $\CN=2^*$ 理论}\label{subsec:N2-star-moduli}

最后，考虑第 5.4 节% TODO \ref{subsec:Formal-N2-star}
讨论过的 $\CN=2^*$ 理论。如前所述，$C$ 是带一个穿刺点的环面。$C$ 的基本群有生成元 $\alpha,\beta,m$，并有一个关系 $\alpha\beta\alpha^{-1}\beta^{-1}=m$。因此，绕穿刺点的单值取固定共轭类的平坦 $SL(2,\IC)$ 联络的模空间 $\CM_{mx}$ 可以这样描述：取满足
\begin{equation}\label{eq:monod-equa-nstar}
ABA^{-1} B^{-1} = M
\end{equation}
的 $SL(2,\IC)$ 矩阵 $A,B,M$，其中 $\Tr M=\mu+1/\mu$ 固定，并对整体的同时共轭作商。

从规范理论的角度看，$\CM_{mx}$ 上的全纯函数代数显然由和乐的迹生成，而这些迹又是 $A^{\pm1}$ 与 $B^{\pm1}$ 所成词的迹。我们断言：全纯函数代数实际上由 $a,b,c$ 生成，其中
\begin{equation}
a:= \Tr A, \qquad  b:= \Tr B, \qquad c:= \Tr AB,
\end{equation}
而 $\CM_{mx}$ 就是满足下式的 $(a,b,c)\in\IC^3$ 构成的空间——这个结果可以追溯到 Fricke：
\begin{equation}\label{eq:moduli-equation}
a^2 + b^2 + c^2 - abc = \mu + 2 + \frac{1}{\mu},
\end{equation}
关于这个模空间的有用信息可参见 \cite{MR2026539,MR1407481}。只要 $\mu\ne -1$，\eqref{eq:moduli-equation} 的解空间就是光滑的。

为论证这些断言，首先注意：由于矩阵是 $2\times2$ 的，只需考虑基础表示中和乐的迹。其次，回忆：若 $x$ 是行列式为 $1$ 的 $2\times2$ 矩阵，则
\begin{equation}\label{eq:CH-thm}
x^2 - \Tr(x) x +1 =0.
\end{equation}
对 $x=A^{\pm 1},B^{\pm 1}$ 使用它，显然可以把 $n_i,m_i\in\IZ$ 的表达式 $\Tr(A^{n_1}B^{m_1}A^{n_2}B^{m_2}\cdots)$ 归约为迹的多项式，其中 $n_i,m_i\in\{0,\pm 1\}$。然后，把 \eqref{eq:CH-thm} 写成 $x^{-1}=\Tr(x)-x$ 的形式来使用，就能用 $A$ 替换 $A^{-1}$、用 $B$ 替换 $B^{-1}$（多项式会随之改变）；最后，对 $x=AB$ 应用 \eqref{eq:CH-thm}，可以把 $\Tr[(AB)^n]$ 归约为 $\Tr(AB)$ 的多项式。此外，利用这些关系可以写出
\begin{equation}
\begin{split}
\Tr M & = \Tr(ABA^{-1} B^{-1} ) \\
& = (\Tr A)^2 - \Tr(ABAB^{-1}) \\
& = (\Tr A)^2 - \Tr(ABA) \Tr(B) + \Tr(ABAB)  \\
& = (\Tr A)^2 + (\Tr B)^2 + (\Tr AB)^2 - (\Tr A )(\Tr AB) (\Tr B) - 2
\end{split}
\end{equation}
由此得到 \eqref{eq:moduli-equation}。

为了重现预期的覆盖模空间
\begin{equation}\label{eq:A1-covering-maps-nstar}
\begin{matrix}
    &  &  \CM_{mx} &  &   \\
    & \swarrow & \downarrow  &  \searrow & \\
    \CM_{1} &  & \CM_{2} &  & \CM_{3} \\
    & \searrow & \downarrow  &  \swarrow & \\
    &  &  \CM &  &   \\
\end{matrix}
\end{equation}
我们注意到：由 \eqref{eq:monod-equa-nstar} 定义的 $\CM_{mx}$ 具有 $\IZ_2\times\IZ_2$ 对称群，由 $(A,B)\to(-A,B)$ 与 $(A,B)\to(A,-B)$ 生成，因而三个非平凡元素对 $(a,b,c)$ 的作用为
\begin{equation}
\begin{split}
g_1: (a,b,c) & \rightarrow (-a,b,-c) \\
g_2: (a,b,c) & \rightarrow (a,-b,-c) \\
g_3: (a,b,c) & \rightarrow (-a,-b,c)\\
\end{split}
\end{equation}
显然 $\CM_i=\CM_{mx}/\langle g_i\rangle$，而 $\CM=\CM_{mx}/(\IZ_2\times\IZ_2)$。注意 $\CM$ 是绕穿刺点单值取固定共轭类的平坦 $PSL(2,\IC)$ 联络的模空间。

相应地，$\IZ_2\times\IZ_2$ 的各元素在 $\CM_{mx}$ 上有六个不动点，例如 ${\rm Fix}(g_3)$ 是 $(0,0,\pm(\mu^{1/2}+\mu^{-1/2}))$。若对其中一个 $\IZ_2$ 子群作商，则有两个 $A_1$ 奇点，它们被剩余的 $\IZ_2$ 对称群互换。例如 $\CM_3$ 在 $[0,0,\pm(\mu^{1/2}+\mu^{-1/2})]$ 处有两个 $A_1$ 奇点，被 $g_1$（或 $g_2$）互换。若对 $\IZ_2\times\IZ_2$ 作商得到 $\CM$，则有 $3$ 个互不相同的 $A_1$ 奇点。在纤维化 $\CM\to\CB$ 下，它们投影为第 5.4 节% TODO \ref{subsec:Formal-N2-star}
提到的 $u$ 平面上的三个奇点。

第 5.4 节% TODO \ref{subsec:Formal-N2-star}
的形式 Wilson 算子与函数 $a,b,c$ 漂亮地对应。确切地说，$y\to 1$ 极限给出
\begin{equation}
\begin{split}
W_1 & \to a, \\
W_2 & \to b, \\
W_3 & \to c.
\end{split}
\end{equation}
明白了这一点，我们就认出方程 (5.28)% TODO \eqref{eq:quant-skein-rel}
正是对应于下式的量子 skein 关系：
\begin{equation}
{\rm Tr}(A) {\rm Tr}(B) = {\rm Tr}(AB) + {\rm Tr}(AB^{-1}).
\end{equation}

注意，这确实是与 Chern--Simons 理论中熟悉的那种形式相同的 skein 关系，因为 $A$、$B$ 两个圈在环面上相交，可以有两种方式解开，系数分别为 $y^{\pm 1/2}$。我们将在第 13 节% TODO \ref{sec:Open-Problems}
再次评论这一点。类似地，我们也能认出方程 (5.29)% TODO \eqref{eq:quant-mod-eq}
是 \eqref{eq:moduli-equation} 的量子类比。（可以验证，当 $y\to 1$ 时 $\sum a_{ijk}\to 1$。）

最后注意，$SL(2,\IZ)$ 的 S-对偶对称在这个模空间中表现得非常漂亮，因为可以把定义方程写成
\begin{equation}
A B = M B A,
\end{equation}
它也可以等价地写成
\begin{equation}
A (B A) = M (B A) A
\end{equation}
\begin{equation}
B A^{-1} = (A^{-1} M A) A^{-1} B.
\end{equation}
因此，我们可以定义 $SL(2,\IZ)$ 的生成元在和乐上的作用：$T: A \to A, B\to B A, M\to M$ 以及 $S: A \to B, B\to A^{-1}, M\to A^{-1} M A$。可以验证 $S^4$ 是恒等变换，但关系 $(ST)^2= S^2$ 在和乐上并不满足。尽管如此，
\begin{equation}
STS\cdot (A,B,M) = B^{-1}\left( T^{-1} S T^{-1} \cdot (A,B,M) \right) B
\end{equation}
因此它在和乐的迹上的作用确实通过模群分解。
